\subsection{Elimination of the binary obstruction and surjectivity of descent}

Continue with the actual discrete groups $B_0,U_0,J_0,L$
of B.2.15 and their specified completion maps.  Thus
$N=L^6\triangleleft U_0$, $U_0/N=E$, and
$J_0/N=J_{\mathcal T}$.  The two groups
$T_2,T_3\triangleleft B_0$ are the normal closures of
$\tau^2,\tau^3$.  In this subsection put
\[
 V=H_1(L,\mathbb F_2),\qquad
 W_2=V_{T_2},\quad W_3=V_{T_3},\qquad
 A=W_2\oplus W_3.                                      \tag{B.S61}
\]
The subscript denotes coinvariants for the actual
conjugation action.  Their dimensions are $25,2,3,5$,
respectively.  The actual quotient map $V\to A$ is onto.
Define
\[
 P=([T_2,L]L^2)\cap([T_3,L]L^2),\qquad
 Q=B_0/P,\qquad R=U_0/P^6.                              \tag{B.S62}
\]
Both factors in the intersection lie in $L$ and are
normal in $B_0$.  These definitions give the exact
commutative diagram
\[
\begin{array}{ccccccccc}
 1&\longrightarrow&L^6&\longrightarrow&U_0&\longrightarrow&E&\longrightarrow&1\\
 &&\downarrow&&\downarrow&&\Vert\\
 1&\longrightarrow&A^6&\longrightarrow&R&\longrightarrow&E&\longrightarrow&1,
\end{array}                                                   \tag{B.S63}
\]
with $|Q|=144\cdot2^5=4608$ and
$|R|=|E|2^{30}=2^{44}3^6$.  In particular all maps
used below are the literal six-coordinate quotient maps.

\begin{result}{B.2.24}{Proved}{The complete reduced central kernel vanishes}
For the universal central extension and the two actual
torsion correction maps of B.2.16 one has
\[
 M=\operatorname{im}\kappa_s\,
                  \operatorname{im}\kappa_r=S,
                  \qquad M/S=1.                           \tag{B.S64}
\]
Consequently the literal descent homomorphism is onto
both its multiplier-one subgroup and its full domain:
\[
                D_P(H)=H,\qquad
 D_P(\widehat{GT}_0)=\widehat{GT}_0.                        \tag{B.S65}
\]
\end{result}

\begin{proof}
We first show that a central character which kills the
two torsion correction images has zero restriction to
$N$.  Let $\chi:M/S\to\mathbb F_2$ be a homomorphism,
and also write $\chi$ for its composite with $M\to M/S$.
Push out $U^{\mathrm c}\to U$ along this homomorphism
and pull the resulting central extension back along
$U_0\to U$.  Let
\[
                     e_\chi\in H^2(U_0,\mathbb F_2)          \tag{B.S66}
\]
be its class.  The named lifts of $s,r$ have orders two
and three.  Its pullback to $B_0$ is split by the actual
map from $B_0$ to $U^{\mathrm c}$.

The restriction to $N$ is an $E$-invariant class.
Since $L$ is free, its only degree-two components are
the pair forms
\[
 H^2(N,\mathbb F_2)=
       \bigoplus_{i<j}(V_i\otimes V_j)^* .                   \tag{B.S67}
\]
The nine pairs with trivial common quotient contribute
zero invariants, as in B.2.21: each has the independent
$B_0\times B_0$ action, and $V_{B_0}=0$ follows from
$[B_0,L]=L$.  The two exceptional matching orbits have
coinvariant dimensions one and two, respectively.
We now retain the actual torsion elements which test
all three of these classes.

For the pair $(1,2)$ the source element $s$ interchanges
the two coordinates without an inner change.  For any
$a$ in the first copy of $L$, put $b=sas^{-1}$ in the
second.  The elements $a,b$ commute and
$c=ab$ belongs to $C_{J_0}(s)$.  In a central extension
write $S$ for the named order-two lift of $s$ and choose
a lift $A_0$ of $a$.  Then $S A_0 S^{-1}$ lifts $b$,
and the commutator of $A_0(S A_0 S^{-1})$ with $S$
is, up to inverse, the central commutator
$[A_0,S A_0 S^{-1}]$.  Its value is the pairing of
$e_\chi|_N$ with $v\otimes v$, where $v$ is the
image of $a$ in $W_2$.  It vanishes because
$\chi\kappa_s(c)=0$.

For the pair $(1,5)$ use the actual involution
$b_0=rsr^{-1}$.  Composing the coordinate actions of
$r$ and $s$ from B.2.13 gives exactly
\[
 \operatorname{Ad}_{b_0}(w_1,w_5)=
       (\tau^{-1}w_5\tau,\ \tau w_1\tau^{-1}).              \tag{B.S68}
\]
For $a$ in the first factor, $a(b_0ab_0^{-1})$
centralizes $b_0$.  Conjugation by $r^{-1}$ sends this
centralizer to $C_{J_0}(s)$ and leaves every central
kernel value unchanged.  The same commutator calculation
therefore makes the pairing with
$v\otimes\operatorname{Ad}_\tau(v)$ zero, for every
$v\in W_3$.

Here is the complete finite-field spanning check for
these actual diagonals.  In the right-module bases
obtained from the integral Bass--Serre graph of B.2.21,
the matrices of $s,r$ are
\[
\begin{array}{c|c|c}
 &M_s&M_r\\ \hline
 W_2&\left(\begin{smallmatrix}1&0\\1&1\end{smallmatrix}\right)
     &\left(\begin{smallmatrix}0&1\\1&1\end{smallmatrix}\right)\\[2mm]
 W_3&\left(\begin{smallmatrix}1&0&0\\1&1&0\\1&0&1\end{smallmatrix}\right)
     &\left(\begin{smallmatrix}1&0&0\\1&1&1\\0&1&0\end{smallmatrix}\right).
\end{array}                                                   \tag{B.S69}
\]
Right multiplication by $M_g$ represents conjugation
by $g^{-1}$.  Thus the second diagonal in (B.S68) is
$v\otimes vM_r^2M_s$, where
\[
       M_r^2M_s=
       \begin{pmatrix}1&0&0\\1&0&1\\1&1&1\end{pmatrix}
                         \quad\hbox{on }W_3.                \tag{B.S70}
\]
For $W_2\otimes W_2$, the relation matrices are
the two actual pair actions
$(M_s,M_sM_r)$ and $(M_sM_r,M_s)$ and the signed
swap, each minus the identity.  Their rank is three
in dimension four; adjoining the four rows
$v\otimes v$ raises the rank to four.
For $W_3\otimes W_3$, the pair actions are
$(M_s,M_s)$ and $(M_sM_r,M_sM_r^2M_s)$; the matching
swap is $-(1\otimes M_s)\operatorname{Swap}$.
Their relation rank is seven in dimension nine;
adjoining the eight rows $v\otimes vM_r^2M_s$
raises the rank to nine.  Thus these torsion diagonals
span both entire pair coinvariant spaces.  The code
\path{work/centralizer_code/check_binary_torsion_diagonals.py}
reconstructs the quotient matrices from all 144
integral graph actions, checks their intertwining
maps, and verifies both ranks by two elimination
implementations.  Its JSON file retains every relation
and diagonal row.  We have proved
\[
                  e_\chi|_{N}=0.                            \tag{B.S71}
\]

Next compare the first filtration steps of second
cohomology in (B.S63).  The maps on the first-homology
rows are induced by the actual dual inclusion
$A^*\hookrightarrow V^*$.  Shapiro identifies the
degree-zero and degree-one maps in that row with
\[
\begin{aligned}
 H^0(J_{\mathcal T},A^*)&\longrightarrow
                         H^0(J_{\mathcal T},V^*),\qquad 0\longrightarrow0,\\
 H^1(J_{\mathcal T},A^*)&\longrightarrow
                         H^1(J_{\mathcal T},V^*),\qquad
                         \mathbb F_2\xrightarrow{\ \sim\ }\mathbb F_2.
\end{aligned}                                                   \tag{B.S72}
\]
These statements concern the indicated maps, not only
equal dimensions.  The exact finite presentation and
graph matrices give, on $A^*$, 90 cocycle unknowns,
relation rank 84, and coboundary rank 5.  On $V^*$ the
corresponding counts are 450, 424, and 25.  The
inclusion of the cocycles has rank one after quotienting
by the 25-dimensional coboundary space.  The two
invariant spaces have dimension zero.  The complete
map is reproduced in
\path{work/centralizer_code/check_binary_first_filtration.py}
and its exported $25\times5$ quotient matrix.

In total degree two, the first cohomology filtration
has only bidegrees $(2,0)$ and $(1,1)$.  For $(2,0)$
the only incoming differential is from $(0,1)$, and
both these source spaces are zero by (B.S72).  The
maps on the zero-th row are identities.  For $(1,1)$
there are no incoming differentials, and its only
outgoing differential is to $(3,0)$; its source map is
the isomorphism in (B.S72), and its target map is the
identity.  Naturality therefore identifies its kernel.
There are no later differentials in either bidegree.
The two finite filtrations give the actual isomorphism
\[
\begin{aligned}
 &\ker\bigl(H^2(R,\mathbb F_2)\to H^2(A^6,\mathbb F_2)\bigr)\\
 &\hspace{2em}\xrightarrow{\ \sim\ }
 \ker\bigl(H^2(U_0,\mathbb F_2)\to H^2(N,\mathbb F_2)\bigr).
\end{aligned}
                                                               \tag{B.S73}
\]
This argument uses only the two maps in (B.S72).
No assertion about degree-two module cohomology is
needed.

By (B.S71)--(B.S73) there is a unique class
$e_R\in H^2(R,\mathbb F_2)$ which inflates to
$e_\chi$ and restricts to zero on $A^6$.
Its pullback to $B_0$ is zero because that is the
pullback of $e_\chi$.  It therefore belongs to the
marked central subspace: the named $s,r$ admit lifts
of their stated orders.

We retain the complete finite calculation of this
marked subspace.  The group $Q$ in (B.S62) is built as
the joint marked quotient of $\mathcal T$, the
mod-two homology cover of $B_0/T_2$, and the quotient
$B_0/[T_3,L]L^2$.  The last quotient has order
$144\cdot2^3=1152$.  Its construction uses precisely
the module $W_3$ in (B.S61): the named generators
have orders two and three and generate the whole
extension, its kernel over $\mathcal T$ has order
eight, and the image of $T_3$ centralizes that kernel.
The resulting map from $B_0/[T_3,L]L^2$ is onto between
groups of the same finite order.  For the first
cover, $L\to H_1(T_2,\mathbb F_2)$ is onto because
$(T_2/L)^{\mathrm{ab}}=C_3$, and both its source
coinvariant quotient and target have dimension two.
This verifies the first factor of the map in (B.S61).
The joint layer has order $2^5$, as checked directly.
Thus the computed marked $Q$ is exactly (B.S62).

The induced six-coordinate source is $R$ in (B.S63).
Its second cohomology has dimension five, and the
restriction to the two source torsion factors has rank
one.  The universal marked elementary central cover
therefore has kernel $C_2^4$.  In its exported four
central coordinates, the complete word calculations
give
\[
\begin{array}{c|c}
 \text{central values}&\text{row-space description}\\ \hline
 \text{within factors}
       &\langle(0,0,1,1)\rangle\\
 \text{between factors}
       &\langle(1,0,1,1)\rangle\\
 \operatorname{im}\kappa_s
       &\langle(1,1,1,0),(0,1,0,1),(1,0,1,0)\rangle.
\end{array}                                                     \tag{B.S74}
\]
Here ``within factors'' denotes the squares and
commutators in one of the six individual factors;
``between factors'' denotes commutators from distinct
factors.  In particular the first two rows have joint rank two,
the last row has rank three, and all rows together
have rank four.

These are full subspaces.  The calculation tests the
squares and all pair commutators of a basis of each
of the six literal order-$2^5$ factors of $A^6$;
it also tests every generator of the full pure
centralizer of $s$ in $R$.  Squares and commutators
determine a central $C_2$ extension of an elementary
abelian two-group: their simultaneous vanishing gives
an elementary abelian extension and hence a split
vector-space sequence.  Thus the first two rows of
(B.S74) give precisely the restriction to $A^6$.
The saved power--conjugate presentation, complete
cocycle basis, and the exact $s,r$ words are replayed
in \path{work/centralizer_code/check_binary_pair_restrictions.g}.
This script verifies the actual source map to $E$,
its elementary kernel of order $2^{30}$, all six
coordinate homomorphisms, the literal six factors,
the generating torsion lifts, and every rank in
(B.S74).  It reads the exported data in
\path{work/centralizer_code/small_homology_cover_2_full.json};
the outputs
\path{work/centralizer_code/check_binary_pair_restrictions.json}
and \path{work/centralizer_code/check_binary_kappa_filtration.json}
retain every central value.  This replay does not
recompute or infer the cohomology from a proper
subgroup of $R$.

The class $e_R$ also annihilates the last row of
(B.S74).  To see the exact centralizer comparison,
the map $J_0\to J_R$ has kernel $P^6$, and $s$ acts
on its factors as three two-cycles, with the factor
maps around each cycle composing to the identity.
Take $h_R\in C_{J_R}(s)$ and lift it to $h\in J_0$.
Then $hsh^{-1}$ and $s$ are order-two elements with
the same image in $R$.  Their difference is a
one-cocycle of the order-two action on $P^6$.
On each permuted pair this cocycle is a coboundary:
choose the first correcting coordinate and solve for
the second; the second equation is the order-two
cocycle identity.  Correcting $h$ in $P^6$ gives
a lift in $C_{J_0}(s)$.  Thus
\[
                       C_{J_0}(s)\longrightarrow C_{J_R}(s)
                     \quad\hbox{is onto}.                    \tag{B.S75}
\]
The equality of extension classes under pullback
intertwines the commutator values for these actual
elements.  Hence $e_R$ evaluates to zero on the entire
finite $\kappa_s$ image, because $\chi$ kills its
upstairs image.  We already know that $e_R$ has zero
restriction to $A^6$.  The rank-four assertion in
(B.S74) now forces $e_R=0$ in the marked central
subspace.  It follows that $e_\chi=0$.

Finally this vanishing implies $\chi=0$ on the actual
profinite $M$.  In the discrete source extension
$K_0\to B_0\to U_0$, the transgression
\[
 \operatorname{Hom}(K_0/[B_0,K_0],\mathbb F_2)
                     \longrightarrow H^2(U_0,\mathbb F_2)
                                                               \tag{B.S76}
\]
is injective: the preceding restriction from
$H^1(B_0,\mathbb F_2)$ is zero, since
$K_0\subseteq B_0'$.  The pullback of $\chi$ along
$K_0\to M$ maps under (B.S76) to $e_\chi$, so it is
zero.  The image of $K_0$ is dense in $M$ by B.2.15,
and $M$ is finite.  Thus $\chi=0$.
The finite group $M/S$ is a two-group by B.2.23.
Every nontrivial finite abelian two-group has a
nonzero character to $\mathbb F_2$, so the result
just proved gives $M/S=1$.  This proves (B.S64).

The exact image identity B.2.19 now gives
$D_P(H)=H$.  For a full target $g\in\widehat{GT}_0$,
choose an actual arithmetic element $\gamma\in G_{\mathbb Q}$
with $\lambda(\gamma)=\lambda(g)$, using the surjective
cyclotomic character.  Then $h=g\gamma^{-1}\in H$.
Choose $a\in H$ with $D_P(a)=h$.  The actual arithmetic
comparison in B.1 gives
$D_P(c\gamma c)=\gamma$.  Hence
\[
                        D_P(a\,c\gamma c)=g.                \tag{B.S77}
\]
This proves full surjectivity with the same literal
descent map, and completes (B.S65).
\end{proof}
